%% APPENDIX E -- radio-frequency interference screening.  Carries
%% app:rfi (05_results.tex, 06_discussion.tex, app_J, counts_r15.py).
%% r30 (write-apx): Table E1 (tab:rfitests) is one paragraph giving each test
%% with its outcome; every value is a macro the table printed.  The rows'
%% populations are named where the reader needs them.  \AppMTubeP carries
%% its own relation sign ($<3\times10^{-4}$), hence "$p$\,\AppMTubeP"; r35: the same holds
%% for \NtTubeCenP, the census-only tube count (N3 misc_r35 S4) that replaces it here.
%% r34 (W3): readability only, one idea per paragraph; no wording of a test changed.
%% r35 (W3, m1): each population is named as a set the paper defines (census,
%% archival tail, hold-out); the frequency-tube test now runs over the census's
%% \CsNCross crossings (\NtTubeCen*); \AppMTube*/\AppMNCrossF (released catalogue rows) no longer cited.
\section{Radio-frequency interference screening}
\label{app:rfi}

Interference is the first alternative explanation for an unattributed
crossing, although an interferometer is a poor receiver of it, since a
signal not in the far field at the phase centre is suppressed by its rapidly
varying geometric phase. Allocations exist across the searched bands
\citep{ITU2024}, and the one documented hazard, spaceborne cloud radar at
\RadarLo--\RadarHi\,GHz \citep{CloudSat2002,EarthCARE2023,ALMAMemo504},
overlaps \RadarNWin{} of the \RadarNTot{} census windows, with \RadarNCross{}
crossings.

The correlator flagging was not retained, so no per-channel interference
flag exists, and interference can be looked for only in the way the
crossings are distributed. The tests below therefore describe the
population. They clear no single crossing, and they do not substitute for
the ON/OFF cadence of a dedicated search. Each test could have failed, and
none did.

Each $p$ value is computed against the same peaks redrawn on their own
channel grids. Crossings cluster in sky frequency only on the survey's own
lines. \NtTubeCenObs{} of the \CsNCross{} census crossings lie within
$\pm\DispoHalfKms$\,km\,s$^{-1}$, in sky frequency, of one of the bright
lines the pipeline searched for from the start, those of CO and its
isotopologues, HCN, HCO$^{+}$, CS, CN, SiO and H$_2$CO, against
\NtTubeCenNull{} expected ($p$\,\NtTubeCenP). That is a coarser test than attribution, which uses the
full list of Appendix~\ref{app:linemask} in the star's own frame, so the
two counts differ.

The \QcNQP{} unattributed quality-passing crossings of the \NtSetUnName{} do not crowd on one
another between \QcTolLo{} and \QcTolHi\,MHz ($p=\QcPQPCorr$), whereas the
attributed ones do. None of the
\TlNCross{} crossings of \TlCrossPop{} lies within \ScRadMHz\,MHz of a
telluric line. Peaks are spread uniformly within their windows (median
\RfiFracMed), which a fixed intermediate frequency would not allow.

Interference reaching several stars would make them cross at the same
frequency. Coincidences between different stars within \RfiOccTol\,MHz number
\RfiOccObs{} against \RfiOccNull{} expected ($p=\RfiOccP$). Of the
\ImNScreen{} unattributed crossings in catalogued blocks, \ImNNoParent{} have
no catalogued transition in the image band and \ImNParentAbsent{} have one that
their own block saw too faint to leak; for \ImNLeft{} the image frequency was not
observed.

The one evident concentration is at the carbon monoxide
tunings. The \TlBandAll{} crossings between \TlBandLo{} and \TlBandHi\,GHz
include \TlBandObs{} of the \TlBandNSub{} unattributed crossings of the
census and the hold-out, against \TlBandNull{} expected from where the
windows were tuned.
